\section{\bt{Discussion and Conclusion}{讨论与结论}}
\label{sec:limitations}
\bi{\system\ shares executable metric and feature definitions, their data conditions, and maintained evidence to reduce repeated database exploration.}
{\system{}共享可执行的指标与特征定义、数据条件及持续维护的证据，减少重复的数据库探索。}

\bi{The controls show reduced duplicate validation, and SQL/version caching achieves similar reuse of identical checks. The snapshot control isolates the effect of retaining the validation snapshot until query execution. The feature-access pilot establishes neither short-horizon amortization nor a composition advantage; content-matched repeated trials remain necessary.}
{对照实验显示重复验证减少，SQL／版本缓存也能实现相近的相同检查复用。快照对照实验分离检验了将验证快照保留至查询执行的作用。特征访问预实验尚未证明短期成本摊销或组合优势，仍需特征内容一致的重复试验。}

\bi{\emph{Future work.} We will support repeated-role and many-to-many joins, incremental feature maintenance, and explicit units, authoritative sources, filter semantics, and value-distribution checks.}
{\emph{未来工作。}我们将支持同表多角色与多对多连接、特征增量维护，并扩展显式量纲、权威来源、过滤语义及取值分布检查。}

\section*{\bt{Acknowledgment of AI-Generated Content}{AI 生成内容声明}}
% ICDE 2027 requires this section (not counted in the 12 pages). The authors must confirm the wording below before submission.
% Disclosure scope follows https://icde2027.github.io/cf-research-papers (checked 2026-10-23).
\bi{Evaluated agents and memory operations use \llmmodel{} under the documented configurations. Under the authors' direction, Claude and ChatGPT assisted with coding, data analysis, visualization, and manuscript editing. The authors designed the study, verified all code, results, and text, and take full responsibility for the work.\looseness=-1}
{受评估的智能体与记忆操作使用 \llmmodel{}，配置见实验说明。在作者指导下，Claude（Anthropic）与 OpenAI Codex 辅助了编码、数据分析、可视化及文稿编辑。作者设计研究、核实全部代码、结果和文本，并对全部工作负责。}
